% HD-F1-T34. GPT-6 Astra Pro, /root/transport; 22 November 2026 UTC.
\section[HD-F1-T34: unitary level-one semilinear exclusion]{HD-F1-T34: exact semilinear descent excludes the remaining unitary level-one sources over fields of size two and three}
\label{sec:f1-unitary-level-one-semilinear-exclusion}

Every remaining T2-OTHER source with projective socle
\(PSU_N(q)\), \(q\in\{2,3\}\), \(N\ge24\), and
level-one constituents on its actual perfect central
layer is excluded. All original restriction
multiplicities, scalar phases, and outer actions are
retained. Combined with T32 and T33, this gives the
effective necessary bounds
\[
 \boxed{\qquad q\le27,\qquad N<10^{62}\qquad}
 \tag{F1T34.1}
\]
for every remaining almost simple linear or unitary
socle \(PSL_N(q)\) or \(PSU_N(q)\). The conclusion
restricts the original source gate; it does not assert
the existence of a source in this finite range.

\paragraph{1. Actual scalar projectors and their field action.}
Put \(A=\mathbf F_{q^2}^{N}\), with
\(\overline{x}=x^q\) and standard Hermitian form
\(h(x,y)=\overline{x}^{,t}y\). Write
\(H=SU_N(q)\) and \(\widehat H=GU_N(q)\).
The genuine total Weil representation \(\mathcal W\)
of \(\widehat H\) has character
\[
 \zeta_N(a)=(-1)^N(-q)^{d(a)},\qquad
 d(a)=\dim_{\mathbf F_{q^2}}\ker(a-1).
 \tag{F1T34.2}
\]
For \(C=\mu_{q+1}\subset\mathbf F_{q^2}^{\times}\)
and \(\eta\in\widehat C\), the actual central
idempotent and its image are
\[
 p_\eta=\frac1{q+1}\sum_{s\in C}
             \eta(s)^{-1}\rho_{\mathcal W}(sI),
 \qquad W_\eta=p_\eta\mathcal W,
 \qquad
 \chi_\eta(a)=\frac1{q+1}\sum_{s\in C}
                 \eta(s)^{-1}(-1)^N(-q)^{d(sa)}.
 \tag{F1T34.3}
\]
The existence of the genuine representation with
character (F1T34.2) is the named published Weil
representation input of Guralnick--Larsen--Tiep,
\emph{Forum of Mathematics, Pi} \textbf{8} (2020),
e2, equation~(1.1) and Example~3.12. The actual
projector is its \(j=1\) instance of equation~(8.8).
Evaluating (F1T34.3) at one gives
\[
 \begin{aligned}
 d_1:=\chi_1(1)&=\frac{q^N+(-1)^Nq}{q+1},\\
 d_\eta:=\chi_\eta(1)&=\frac{q^N-(-1)^N}{q+1}
                         \quad(\eta\ne1),\\
 d_\eta&\ge\frac{q^N-q}{q+1}.
 \end{aligned}
 \tag{F1T34.4}
\]
Here is a direct proof of all required sector
irreducibility assertions. The group \(H\) has
exactly \(q+1\) orbits on \(A\): zero, nonzero
isotropic vectors, and one orbit for each nonzero
value of \(h(v,v)\). For nonisotropic vectors of
equal norm, extend each vector by an orthogonal
Hermitian basis. The complements are isometric:
successive orthogonal splitting gives diagonal
forms, and the surjectivity of
\(\mathbf F_{q^2}^{\times}\to\mathbf F_q^{\times}\),
\(x\mapsto x^{q+1}\), makes all nonzero diagonal
entries equivalent. For an isotropic nonzero
\(v\), choose \(w\) with \(h(v,w)=1\), and replace
\(w\) by \(w-tv\), where
\(t+\overline t=h(w,w)\). The finite-field trace
map is surjective, so this constructs an isotropic
partner and a hyperbolic plane with first vector
exactly \(v\). Complete its orthogonal complement
as before. These constructions give the required
\(\widehat H\)-isometries. Their determinants lie
in \(\mu_{q+1}\); multiply an anisotropic line in
the target orthogonal complement by the inverse
determinant. That complement has dimension \(N-1\)
or \(N-2\), respectively, and is nonzero here.
The correction fixes the target vector, preserves
the Hermitian form, and gives determinant one.
Thus it proves the asserted \(H\)-orbits by actual
maps. The nonzero norm values all occur by norm
surjectivity, and a nonzero isotropic vector is
\(e_1+ae_2\) with \(a\overline a=-1\).

By (F1T34.2), \(|\zeta_N(h)|^2=q^{2d(h)}\) is the
number of natural vectors fixed by \(h\). Averaging
this count over \(H\), or equivalently counting
each orbit by its point stabilizers, gives
\(\langle\zeta_N,\zeta_N\rangle_H=q+1\).
The \(q+1\) actual sectors in (F1T34.3) all have
positive dimension by (F1T34.4). Semisimplicity now
forces their restrictions to be irreducible and
pairwise inequivalent: any further decomposition
or repeated irreducible would raise the sum of
multiplicity squares above \(q+1\). Their dimensions
exceed one in the present range, so all have level
one, and the least-power definition of level
(Definition~2 and Lemma~7.1 of the same paper)
shows that they are every level-one constituent.
In particular there is no lower-level summand to
remove. This directly proves the \(j=1\) assertions
called well known in the published page~67
paragraph; its Corollary~8.6 provides a separate
inspected restriction comparison. The estimates
and semilinear maps below are also proved explicitly.

For any \(a\), applying entrywise conjugation to
the kernel gives
\(d(s\overline a)=d(\overline s a)\).
Since \(\overline s=s^{-1}\) on \(C\), substitution
in (F1T34.3) yields
\[
                  \chi_\eta(\overline a)
                         =\chi_{\eta^{-1}}(a).
 \tag{F1T34.5}
\]
Conjugation by \(\widehat H\) preserves every sector,
whereas the field automorphism can exchange the
sectors \(\eta\) and \(\eta^{-1}\). The original
isotypic summands are therefore not assumed to be
individually fixed by every outer element. An element
exchanging two such summands has trace zero on their
sum, by its off-diagonal block matrix.

\paragraph{2. Direct estimates for inner and diagonal lifts.}
For a nonscalar \(a\in\widehat H\), set
\(d_s=d(sa)\). The distinct natural eigenspaces
have direct sum, so
\[
 0\le d_s\le N-1,\qquad \sum_{s\in C}d_s\le N.
 \tag{F1T34.6}
\]
For integers \(u\ge v\ge1\),
\[
 q^{u+1}+q^{v-1}-q^u-q^v
                 =(q-1)(q^u-q^{v-1})\ge0.
 \tag{F1T34.7}
\]
First add any unused dimension until the sum of the
exponents is \(N\), without exceeding \(N-1\) in
any entry. Then repeatedly apply (F1T34.7). This
concentrates the exponents at \(N-1,1,0,\ldots,0\),
and proves
\[
 \sum_{s\in C}q^{d_s}\le q^{N-1}+2q-1,
 \qquad
 \frac{|\chi_\eta(a)|}{d_\eta}
     \le\frac{q^{N-1}+2q}{q^N-q}\le\frac35
                              \quad(N\ge6).
 \tag{F1T34.8}
\]
For the final inequality, clearing denominators gives
\(q^{N-1}(3q-5)\ge13q\). At \(N=6\), its left
side is \(32\) for \(q=2\), and \(972\) for
\(q=3\), exceeding respectively \(26\) and \(39\);
it increases with \(N\).

Keep the original \(G,V,Z,L\) and the actual
surjection \(\beta:H\to L\) of T32. The universal
central-cover exceptions for this unitary family have
\(N\le6\), so the exact comparison square
\[
             \beta\widetilde a_g=c_g\beta
 \tag{F1T34.9}
\]
applies throughout the present range. The inflated
restriction is the actual direct sum
\[
 V_{\mathbf C}|_H=\bigoplus_\eta(W_\eta\otimes M_\eta),
 \quad M_\eta=\operatorname{Hom}_H(W_\eta,V_{\mathbf C}),
 \quad w\otimes f\longmapsto f(w),
 \quad m=\sum_\eta d_\eta\dim M_\eta.
 \tag{F1T34.10}
\]
Zero multiplicity spaces are omitted. This evaluation
map is an isomorphism by the inspected irreducibility
and ordinary semisimplicity of the actual restriction.

Let the original \(g\) induce conjugation by
\(a\in\widehat H\). It preserves each summand in
(F1T34.10), and its action there is exactly
\[
 A_{g,\eta}\otimes T_{g,\eta},\qquad
 A_{g,\eta}=\rho_{W_\eta}(a),\qquad
 T_{g,\eta}(f)=\rho_V(g)\,f\,A_{g,\eta}^{-1}.
 \tag{F1T34.11}
\]
The square (F1T34.9) verifies directly that
\(T_{g,\eta}(f)\) is \(H\)-linear. Average a positive
Hermitian form over the finite original group. On
each isotypic summand its pullback under evaluation is
the product of a fixed \(\widehat H\)-invariant form
on \(W_\eta\) and a form on \(M_\eta\), by Schur's
lemma. Thus \(T_{g,\eta}\) is unitary, with the
original scalar phase retained. If \(g\) is
nonscalar, then \(a\) is nonscalar, since the
centralizer of the projective socle in the almost
simple projective group is trivial. Summing the
literal traces in (F1T34.11) gives
\[
                   |\operatorname{Tr}_V(g)|\le3m/5.
 \tag{F1T34.12}
\]
For a unitary operator on an \(m\)-space with
\(r=\operatorname{rank}(g-1)\), the fixed eigenvalues
contribute \(m-r\), and the remaining eigenvalues
have total real part at least \(-r\). Hence
\(\operatorname{Re}\operatorname{Tr}(g)\ge m-2r\).
Equation (F1T34.12) proves
\[
                   \operatorname{rank}(g-1)\ge m/5
 \tag{F1T34.13}
\]
for every nonscalar inner or diagonal lift. Every
nonidentity scalar has rank \(m\). Multiplying any
matrix in (F1T34.12) by a scalar of absolute value
one preserves its trace bound.

\paragraph{3. The actual semilinear representative and its square.}
Steinberg, \emph{Lectures on Chevalley Groups},
Theorem~36, and the twisted-case passage following
Malle--Testerman Theorem~24.24 give the exact
inner, diagonal, and field automorphisms. Since
\(q=2,3\) is prime, every remaining outer action
is represented on the displayed natural module by
\[
       \tau=a\circ\overline{\phantom{x}},
       \qquad a\in\widehat H,
       \qquad \tau^2=a\overline a\in H.
 \tag{F1T34.14}
\]
Indeed \(\det(a)\in\mu_{q+1}\), so
\(\det(a\overline a)=1\). For the standard
Hermitian form the inverse-transpose operation on
\(H\) is entrywise conjugation; there is no omitted
additional graph action. The universal-cover square
identifies these representatives with the original
action of \(g\).

When \(\tau^2\) is nonscalar, the original matrices
satisfy
\[
                 g^2=z\,\beta(\tau^2),\qquad z\in Z.
 \tag{F1T34.15}
\]
To see this, their projective conjugations on the
socle agree; their quotient is in its trivial
projective centralizer and therefore lies in \(Z\).
The bound (F1T34.12) applies to the right side,
including its actual scalar phase. Consequently
\[
 \operatorname{rank}(g-1)
 \ge\operatorname{rank}(g^2-1)\ge m/5,
 \tag{F1T34.16}
\]
where the first comparison is the literal inclusion
\(\ker(g-1)\subset\ker(g^2-1)\).
This argument does not require \(g\) to preserve
each summand in (F1T34.10).

\paragraph{4. Scalar-square descent, with the exact form multiplier.}
It remains to consider \(\tau^2=cI\). Commuting
\(\tau\) with its square gives \(\overline c=c\),
so \(c\in\mathbf F_q\). Since its square is unitary,
\(c^2=1\). Thus \(c=1\) for \(q=2\), and
\(c=\pm1\) for \(q=3\).
Choose \(u\in\mathbf F_{q^2}^{\times}\) with
\(u\overline u=c\), and put
\[
 J=u^{-1}\tau,
 \qquad J^2=1,
 \qquad h(Jx,Jy)=c^{-1}\overline{h(x,y)}.
 \tag{F1T34.17}
\]
For \(c=1\), take \(u=1\). For \(q=3,c=-1\),
write \(\mathbf F_9=\mathbf F_3[i]\), \(i^2=-1\),
and take \(u=1+i\), whose norm is \(-1\).
The multiplier \(c^{-1}\) in (F1T34.17) is retained.

Let \(U=\ker(J-1)\), viewed over \(\mathbf F_q\).
For \(\theta\in\mathbf F_{q^2}\setminus\mathbf F_q\)
and \(v\in A\), both
\(v+Jv\) and \(\theta v+\overline\theta Jv\)
belong to \(U\); their coefficient matrix has
determinant \(\overline\theta-\theta\ne0\).
Thus \(U\) spans \(A\) over \(\mathbf F_{q^2}\).
Conversely, a relation between \(\mathbf F_q\)-independent
members of \(U\) has coefficients \(a_i+\theta b_i\).
Subtracting its image under \(J\) forces first all
\(b_i=0\), then all \(a_i=0\). Therefore the actual
scalar-extension map is an isomorphism
\[
        U\otimes_{\mathbf F_q}\mathbf F_{q^2}
                   \xrightarrow{\ \sim\ }A.
 \tag{F1T34.18}
\]

When \(c=1\), \(h|_U\) is a nondegenerate
\(\mathbf F_q\)-valued symmetric form \(B\).
When \(q=3,c=-1\), take \(\kappa=i\); then
\(B=\kappa^{-1}h|_U\) is a nondegenerate
\(\mathbf F_3\)-valued alternating form.
For the latter assertion, (F1T34.17) gives
\(\overline{h(x,y)}=-h(x,y)\) on \(U\), and
Hermitian symmetry gives \(B(y,x)=-B(x,y)\).
Nondegeneracy in both cases follows because \(U\)
spans \(A\). In the alternating case \(N\) is even.

A nondegenerate symmetric form in characteristic two
splits orthogonally into blocks of dimensions one or
two: split off a line whenever \(B(v,v)\ne0\), and
the remaining alternating form splits into
nondegenerate planes by choosing a pair with
nonzero pairing and taking its orthogonal complement.
In odd characteristic a symmetric form splits into
lines by the same procedure, since a symmetric form
with all diagonal pairings zero is zero. The
alternating case splits into planes. We thus obtain
an explicit basis change
\(t:\mathbf F_{q^2}^{N}\to A\) for which the Gram
matrix is \(\kappa B\), with \(\kappa=1\) in the
symmetric case, and
\[
                  \tau=t\,(u\overline{\phantom{x}})\,t^{-1}.
 \tag{F1T34.19}
\]
This is the needed descent of the actual semilinear
operator, not a replacement of its representation.

\paragraph{5. A determinant-one block map and the large commutator.}
List the orthogonal blocks from paragraph~4, with
dimensions \(w_i\in\{1,2\}\). Choose a prefix whose
dimension \(d\) is the last partial sum at most
\(N/2\). Since each next block has dimension at most
two,
\[
                    N/2-2<d\le N/2.
 \tag{F1T34.20}
\]
Let \(\lambda\) generate \(\mu_{q+1}\). Define
\(e_0\) to be scalar \(\lambda\) on the selected
blocks and scalar one on the rest. It preserves
\(\kappa B\) and has determinant \(\lambda^d\).
Correct its determinant by multiplying one block of
dimension \(w\) by \(\lambda^b\), where
\[
                       bw\equiv-d\pmod{q+1}.
 \tag{F1T34.21}
\]
This equation is always solvable in the present
actual block decompositions. For \(q=2\), both
one and two are invertible modulo three. For
\(q=3,c=1\), all blocks have dimension one.
For \(q=3,c=-1\), all blocks have dimension two
and \(d\) is even, so \(2b\equiv-d\pmod4\)
is solvable. The corrected matrix \(e\) is an
actual determinant-one isometry of \(\kappa B\),
and differs from \(e_0\) on at most two dimensions.
Thus \(e'=tet^{-1}\in H\).

Conjugation by (F1T34.19) is entrywise conjugation
on these block matrices. Hence the actual natural
and coefficient commutators satisfy
\[
 h_0=[\tau,e']=t(\overline e e^{-1})t^{-1}\in H,
 \qquad
 [g,\beta(e')]=\beta(h_0).
 \tag{F1T34.22}
\]
The second equality follows directly from
(F1T34.9). Before the one-block correction,
\(\overline e_0 e_0^{-1}\) has eigenvalues
\(\lambda^{q-1}\ne1\) and one, with multiplicities
\(d\) and \(N-d\). Changing at most two diagonal
dimensions makes every eigenspace have dimension
at most \(\max(d,N-d)+2<N/2+4\). In particular
\[
                       \operatorname{supp}(h_0)\ge N/2-4.
 \tag{F1T34.23}
\]
For any natural matrix with largest eigenspace
dimension \(\delta\), the actual average (F1T34.3)
gives \(|\chi_\eta(h_0)|\le q^\delta\).
Together with (F1T34.4) and (F1T34.23), this yields
\[
 \frac{|\chi_\eta(h_0)|}{d_\eta}
 \le\frac{(q+1)q^{-N/2+4}}{1-q^{1-N}}
 <\frac1{32}\qquad(N\ge24).
 \tag{F1T34.24}
\]
Indeed the numerator is at most \(4\cdot2^{-8}\)
and the denominator is greater than \(1/2\).
Summing the actual multiplicities in (F1T34.10)
gives \(|\operatorname{Tr}_V\beta(h_0)|<m/32\),
and therefore
\(\operatorname{rank}(\beta(h_0)-1)>31m/64\).
For \(e_L=\beta(e')\), the exact identity
\[
 [g,e_L]-1
  =\bigl((g-1)e_L-e_L(g-1)\bigr)g^{-1}e_L^{-1}
 \tag{F1T34.25}
\]
has rank at most \(2\operatorname{rank}(g-1)\).
Thus every scalar-square outer lift satisfies
\[
             \operatorname{rank}(g-1)>31m/128>m/5.
 \tag{F1T34.26}
\]
All scalars in the original coefficient matrix
remain in this identity; none has been independently
adjusted in deriving the rank inequality.

\paragraph{6. The original branch contradiction and framework maps.}
The preceding cases exhaust all original nonidentity
inertias. Each has moving rank at least \(m/5\),
including those whose field action exchanges the
actual Weil constituents. For the original
product-one generating tuple, the retained
Chevalley--Weil identity and branch requirement give
\[
      2m+6=\sum_{\nu=1}^{b}\operatorname{rank}(g_\nu-1)
                  \ge\frac{bm}{5}\ge\frac{12m}{5}.
 \tag{F1T34.27}
\]
Hence \(m\le15\), contrary to the original gate
\(m\ge27\). This proves the stated exclusion and,
with T32--T33, the bound (F1T34.1).

Every comparison in this proof is the scalar
projector (F1T34.3), the actual central-cover square,
the evaluation map (F1T34.10), the scalar-extension
isomorphism (F1T34.18), the displayed basis change,
or the literal coefficient commutator. The original
cover and its coefficient local system are unchanged.
The comparisons on middle-extension stalks and
parabolic cohomology are therefore identities, as
are the original cup pairing, rational primitive
idempotent and its integral-cleared Jacobian image,
the actual \(K\)-corner, restricted polarization,
Poincar\'e biextension, Hodge filtration, period
differential, and the Grothendieck--Adams
realizations of F01--F04. Their pairing multiplier
is one. No auxiliary cover or new detector
component is introduced.
